% =====================================================================================
%  Chapter 4 — Control Architectures
%  Thesis: Comparison of Control Strategies for Interaction-Force Aware Multirotor Teams
%
%  DRAFT, 23 August 2026 (writing-track item W.8). Source: docs/18_Strategy_Descriptions.md.
%
%  THE RULE THIS CHAPTER MUST NOT BREAK.
%  Each strategy is presented in two separated parts: what WE implement (verifiable against
%  firmware_app/src/lib.rs and residual_nn.rs) and the REFERENCE METHOD it follows (cited,
%  with deviations stated). These are never merged into one sentence. Our implementation is
%  not the paper's method, and writing as though it were would misrepresent both -- and
%  would undermine RQ1, whose validity rests on each strategy being a defensible
%  representative of its family rather than a reproduction of a specific paper.
%
%  Implementation status is stated honestly per strategy. Several are not implemented.
%  Do not describe an unimplemented strategy as though it exists.
%
%  Citation baseline: 2026-08-27 (see docs/thesis/CITATION_BASELINE.md).
%  After that date: only NEW cites or EDITED prior-work sentences need re-checking.
%
%  2026-09-14: renumbered to match docs/01_Thesis_Project_Snapshot.md's canonical numbering
%  (Strategy 1 = Pure INDI ... Strategy 7 = Geometric + Residual RL). This chapter previously
%  used its own numbering (Strategy 1 = the uncompensated baseline, shifting everything else
%  up by one), inconsistent with docs/01, docs/19 and docs/18. The baseline is now Strategy 0,
%  outside the seven-method count, matching docs/19's usage. Section order changed to match:
%  Strategy 6 (MPC) now follows Strategy 5 (residual RL, cascaded) rather than preceding it,
%  since canonical order is 0,1,2,3,4,5,6,7. Also corrected a factual error introduced in the
%  2026-09-14 substrate-framing removal: Strategy 7 (residual RL on the geometric baseline)
%  DOES use the geometric position/attitude loop -- it was wrongly grouped with the
%  structurally-different strategies.
% =====================================================================================

\chapter{Control Architectures}
\label{ch:control}

Chapter~\ref{ch:model} defined the residual wrench and the two ways of obtaining it. This
chapter describes the control strategies compared in this thesis, and in each case the point
at which knowledge of the interaction enters the control law.

Four controllers: an uncompensated baseline (Strategy~0), flown in every scenario as the
reference condition, and the three strategies compared against it (Strategies~1--3).
Table~\ref{tab:ctrl-summary} lists them, and the rest of the chapter takes them in order.

\begin{table}[t]
  \centering
  \caption{The strategies compared, and the implementation status of each at the time of writing.
    ``Knows'' refers to the residual: whether the controller has a measurement of it, a prediction
    of it, both, or neither.}
  \label{tab:ctrl-summary}
  \small
  \begin{tabular}{@{}clllcc@{}}
    \toprule
    \# & Strategy & Knows & Enters at & Rotor speed & Status \\
    \midrule
    0 & Geometric baseline           & nothing    & ---                        & no  & flying \\
    1 & Pure INDI                    & measured   & $a_{\mathrm{des}}$, torque & yes & flying \\
    2 & Geometric + NN               & predicted  & $a_{\mathrm{des}}$         & no  & implemented, unflown \\
    3 & FBL + NN                     & predicted  & linearising law            & no  & not implemented \\
    \bottomrule
  \end{tabular}
\end{table}

The chapter treats each strategy as its own controller. What makes the comparison fair is not a
claimed shared implementation, but the experimental conditions held constant for every one of them — the same
airframe, trajectories, formations, separations, residual definition and gains frozen per strategy
after tuning, stated precisely in Section~\ref{sec:ctrl-fairness}.

Each strategy is described in two parts, kept deliberately separate: what is implemented here, and
the published method it follows. They follow published \emph{families}, and the deviations from
each source paper are listed in that strategy's section — none is a bit-identical replica.
Strategies whose author code has not yet arrived stay unimplemented rather than guessed.

% -------------------------------------------------------------------------------------
\section{The position and attitude loop}
\label{sec:ctrl-common}

Strategies~0, 1 and~2 (the geometric baseline, pure INDI and geometric+NN) each use a geometric
position loop and an attitude law realising the corresponding thrust direction; Strategy~1 replaces
the attitude law's inner term with an incremental one. This is stated once here, in its own section,
purely to avoid repeating the same equation three times. Strategy~3 uses a different published
structure entirely, feedback linearisation, described in its own section instead.

\subsection{Position loop}

The desired acceleration is
%
\begin{equation}
  a_{\mathrm{des}}
  \;=\; \underbrace{a_{\mathrm{ff}}}_{\text{trajectory}}
      + K_p e_p + K_v e_v + K_i \!\int\! e_p \,\mathrm{d}t
      + g e_3
      \;\underbrace{-\; a_{\mathrm{res}}}_{\text{measured}}
      \;\underbrace{-\; \hat{a}_{\mathrm{res}}}_{\text{predicted}},
  \label{eq:ctrl-ades}
\end{equation}
%
with $e_p = p_{\mathrm{des}} - p$ and $e_v = v_{\mathrm{des}} - v$, and gains diagonal but distinct
in the horizontal and vertical axes. The commanded thrust vector is $m\,a_{\mathrm{des}}$; its
magnitude projected onto the body $z$-axis gives the collective thrust, and its direction, together
with the desired yaw, gives the desired attitude $R_{\mathrm{des}}$ through the usual flatness
construction.

The last two terms of \eqref{eq:ctrl-ades} are the whole subject of this thesis, and both are
subtracted, for the reason derived in Section~\ref{sec:model-sign}. \emph{Which residual
information is present is what distinguishes Strategies 0, 1 and 2 from one another}, and nothing
else about the position loop changes between them: Strategy~0 has neither term, Strategy~1 the
measured one, and Strategy~2 the predicted one. Strategy~3 injects its prediction elsewhere, in
the feedback-linearising law of Section~\ref{sec:ctrl-s3}.

\subsection{Attitude loop}

The attitude error is the standard geometric one,
$e_R = \tfrac{1}{2}\left(R_{\mathrm{des}}^\top R - R^\top R_{\mathrm{des}}\right)^\vee$, with
$e_\omega = \omega - \omega_{\mathrm{des}}$. In the geometric configuration the torque is
%
\begin{equation}
  \tau = -K_R e_R - K_\omega e_\omega + \omega \times J\omega
         - k_{i,\mathrm{att}} \textstyle\int e_R \, \mathrm{d}t,
  \label{eq:ctrl-tau-geo}
\end{equation}
%
the third term compensating the gyroscopic coupling. This is a reduced form of the law
in~\cite{lee2010geometrictracking}: it drops that work's angular-feedforward bracket and adds an
integral term the source does not have --- see \eqref{eq:ctrl-tau-lee} and the discussion in
Section~\ref{sec:ctrl-s0}. The INDI configurations replace
\eqref{eq:ctrl-tau-geo} with an incremental law, described in Section~\ref{sec:ctrl-s1}.

\subsection{Runtime selection}
\label{sec:ctrl-runtime}

The control law is selected by a parameter at runtime rather than by recompiling: one bit enables
the measured residual in the position loop, another selects the incremental attitude law, and a
third parameter enables the learned prediction. All are pushed to the vehicle at connection time
from a single configuration file.

This matters more than it may appear. Comparing strategies that require different firmware
introduces a confound that is impossible to rule out afterwards, a flash cycle changes more than
the control law, and strategies flown on different days meet different battery states, different
temperatures and different motion-capture calibrations. Runtime selection allows the strategies to
be interleaved within a single session, which is what makes the comparison a paired one.

% -------------------------------------------------------------------------------------
\section{Strategy 0: uncompensated geometric baseline}
\label{sec:ctrl-s0}

\paragraph{Implemented.} Equations~\eqref{eq:ctrl-ades} and \eqref{eq:ctrl-tau-geo} with both
residual terms zero. No knowledge of the interaction whatsoever.

\paragraph{Reference.} Geometric tracking on $SE(3)$~\cite{lee2010geometrictracking}, which tracks
a position trajectory together with a single heading direction, four outputs, not a full
six-degree-of-freedom pose. The attitude error of \eqref{eq:ctrl-tau-geo} is that work's
$e_R = \tfrac{1}{2}(R_d^\top R - R^\top R_d)^\vee$, verified against the source rather than
reproduced from memory.

\paragraph{The implemented torque law is a reduced form, and is labelled as such.} The law in the
source is
%
\begin{equation}
  M = -k_R e_R - k_\Omega e_\Omega + \Omega \times J\Omega
      - J\!\left( \hat{\Omega}\, R^\top R_d \Omega_d - R^\top R_d \dot{\Omega}_d \right),
  \label{eq:ctrl-tau-lee}
\end{equation}
%
with $e_\Omega = \Omega - R^\top R_d \Omega_d$. Equation~\eqref{eq:ctrl-tau-geo} retains the
proportional, damping and gyroscopic terms but omits the final feedforward bracket, and adds an
integral term on attitude error that the source does not contain. It is therefore a
\emph{reduced} law in the same family, not the torque law of~\cite{lee2010geometrictracking}, and
no stability result proved there transfers to it unexamined. Restoring the feedforward bracket
would be a firmware change rather than a wording change, and is outside the scope of this draft.

\paragraph{What the source proves.} Not global asymptotic tracking. With the attitude error
function $\Psi$, exponential stability is established on $\Psi(0) < 1$, initial attitude errors
below \SI{90}{\degree}, and almost-global exponential attractiveness on $1 \le \Psi < 2$,
below \SI{180}{\degree}. The distinction matters here only in that this thesis claims neither: the
baseline is flown as a reference condition, not as a certified controller.

\paragraph{Role.} Strategy~0 is the uncompensated baseline of Table~\ref{tab:ctrl-summary}. It is
flown in every scenario so that ``how much did compensation help'' has a denominator, and so that
the interaction can be separated from the model error discussed in
Section~\ref{sec:model-measure-fail}. It is also the controller under
which all training data is collected, precisely because it does not compensate: a residual observed
under a compensating controller is the residual that survived compensation, which is not what a
predictive model should be fitted to.

% -------------------------------------------------------------------------------------
\section{Strategy 1: pure INDI}
\label{sec:ctrl-s1}

\paragraph{Implemented.} The measured residual of \eqref{eq:model-ares} enters the position loop as
$-a_{\mathrm{res}}$ in \eqref{eq:ctrl-ades}. In the attitude loop the geometric torque law is
replaced by an incremental one. A reference angular acceleration is formed from the geometric
errors,
%
\begin{equation}
  \alpha_{\mathrm{ref}} = \alpha_{\mathrm{des}} - K_R e_R - K_\omega e_\omega,
  \label{eq:ctrl-alpharef}
\end{equation}
%
where $\alpha_{\mathrm{des}}$ is the feedforward angular acceleration obtained from trajectory snap
by differential flatness. The achieved angular acceleration $\alpha_{\mathrm{meas}}$ is obtained by
differentiating the gyro signal and low-pass filtering it, and the torque currently being applied,
$\tau_{\mathrm{cur}}$, is reconstructed from measured rotor speeds. The command is then the
increment
%
\begin{equation}
  \tau = \tau_{\mathrm{cur}} + J\left(\alpha_{\mathrm{ref}} - \alpha_{\mathrm{meas}}\right).
  \label{eq:ctrl-indi-tau}
\end{equation}
%
Because \eqref{eq:ctrl-indi-tau} corrects relative to what is actually being achieved, it rejects
any disturbance appearing in $\alpha_{\mathrm{meas}}$ without needing a model of it. The
differentiation and filtering are the source of the delay identified in
Section~\ref{sec:model-measure-fail}; the filter cutoff sets the trade-off between that delay and
the noise admitted from the gyro.

The position and attitude increments are independently selectable, giving four configurations:
geometric, position-INDI only, attitude-INDI only, and full INDI. Strategy~1 refers to the full
configuration, and that is the controller the comparison reports. The intermediate ones are
diagnostic bits used in the two-robot results chapter to attribute an effect to one loop or the
other. They are not additional strategies and will not be reported as Strategy~1.

\paragraph{Reference.} The incremental formulation with flatness-derived
feedforward follows Tal and Karaman~\cite{tal2018indi}; the cascaded arrangement of position and
attitude INDI follows Smeur et al.~\cite{smeur2017cascadedindi}.

\paragraph{Deviations.} The airframe, gains and filter cutoffs are ours and were identified for
this platform. Tal and Karaman obtain rotor speed from optical encoders~\cite{tal2018indi}; this
implementation uses the platform's own rotor-speed measurement. No claim is made that this
implementation reproduces the tracking performance reported in either work, that would require a
like-for-like replication which is not part of this thesis.

\paragraph{Requirements.} Rotor-speed sensing, without which $\tau_{\mathrm{cur}}$ and
$a_{\mathrm{res}}$ are both unavailable and the method degrades silently
(Section~\ref{sec:model-measure-fail}). The method needs no training data; the residual is formed
online from inertial and rotor-speed measurements.

% -------------------------------------------------------------------------------------
\section{Strategy 2: geometric with a learned residual}
\label{sec:ctrl-s2}

\paragraph{Implemented.} The geometric controller of Section~\ref{sec:ctrl-s0} with the deep-sets
prediction $\hat{a}_{\mathrm{res}}$ of \eqref{eq:model-deepsets} subtracted in
\eqref{eq:ctrl-ades}. The network is evaluated onboard at the control rate from the relative states
of the neighbours, which the flight stack already broadcasts, so the input requires no additional
communication. Weights are trained offline on data collected under Strategy~0 and uploaded before
flight.

The prediction is computed and logged \emph{whether it is enabled or not}. This is what makes the
model's own accuracy measurable independently of its effect on the vehicle: on a flight where the
prediction is not driving the controller, predicted and measured residuals are two independent
quantities and their comparison is meaningful. Once the prediction is driving the controller, the
measured residual is the one that survived compensation and the comparison no longer isolates model
quality. Both flights are therefore needed, and Chapter~\ref{ch:setup} schedules them.

\paragraph{Reference.} The approach of predicting the interaction from relative states with a
permutation-invariant network and applying it as feedforward follows the Neural-Swarm
line~\cite{shi2020neuralswarm,shi2022neuralswarm2}.

\paragraph{Deviations.} These are substantial and are stated plainly. The architecture here is a
deep-sets model of our own sizing, \num{987} parameters, rather than a reproduction of any
published network. Neural-Swarm2 applies spectral normalisation to obtain stability
and generalisation guarantees~\cite{shi2022neuralswarm2}; this implementation does not, and the
guarantees that follow from it are correspondingly not claimed. Neural-Swarm2 also uses the learned
residual in motion planning as well as control~\cite{shi2022neuralswarm2}, whereas here the
trajectories are planned in advance and the residual is used for control only. This strategy is a
representative of the predictive family, and should be read as such rather than as an evaluation of
Neural-Swarm2.

\paragraph{Requirements.} Neighbour relative states; training data. No rotor-speed sensing.

% -------------------------------------------------------------------------------------
\section{Strategy 3: feedback linearisation with a flatness-preserving residual}
\label{sec:ctrl-s3}

\paragraph{Implemented.} \textbf{Not implemented.} The controller code was requested from the
authors and had not been received at the time of writing. The strategy does not block the others,
and its absence is a scope limitation rather than a failure of the comparison
(the discussion chapter).

\paragraph{Reference.} Hsieh et al.~\cite{hsieh2026flatness}, building on the theory of Yang et
al.~\cite{yang2025flatnessresiduals}. The distinguishing property is where the prediction enters:
not as a feedforward term added to a desired acceleration, but inside a feedback-linearising law
whose validity depends on the augmented system remaining differentially flat. A generic learned
residual can destroy that flatness; residuals with a lower-triangular structure, a structured
correction inside the state derivative, not a free wrench, preserve it, and with it the original
flat output~\cite{yang2025flatnessresiduals}. That theory is established in simulation on a planar
quadrotor; the hardware demonstration is Hsieh et al.'s.

Its relevance to this comparison is that it is the strongest published counter-argument to treating
``predictive compensation'' as one method. Hsieh et al.'s abstract reports a \SI{31}{\percent}
reduction in average tracking error; on hardware the body reports \SI{24}{\percent} RMSE for
ROM-FBL over nominal FBL and a further \SI{29}{\percent} for learned FBL over ROM-FBL, matching
the tracking error published by Chee et al.\ for nonlinear model predictive control, at an order
of magnitude less computation in simulation, from under \SI{30}{\second} of training data, and
at a \SI{5}{\milli\second} loop rate obtained \emph{off-board}~\cite{hsieh2026flatness}. Those are
the figures a predictive strategy in this class should be measured against, with the caveat that
none of them was produced by a controller running on the vehicle.

\paragraph{Requirements.} Neighbour relative states; training data. No rotor-speed sensing.

% -------------------------------------------------------------------------------------
\section{What makes the comparison fair, and what it does not guarantee}
\label{sec:ctrl-fairness}



\paragraph{What is held constant.} The airframe, the trajectories, the formations, the
separations, the residual definition and the estimator. Each strategy is tuned for itself; those
gains are then frozen before the comparison flights. This holds uniformly for every strategy — none
is exempted or held to a looser standard on the tasks or the measurement pipeline.

\paragraph{What this does not guarantee.} Two limitations are worth stating before results are
presented rather than after.

Every strategy follows a published \emph{family}, not a specific paper reproduced exactly; the
deviations from each source are listed in that strategy's own section. Strategies whose author code
has not yet arrived stay unimplemented rather than guessed. Because each strategy is its own
controller, ``identical conditions'' means identical trajectories, formations, hardware and
separations — not identical control structure or one gain set imposed on every law. A difference
in outcome could in principle be attributable to a
controller's own structure rather than to how it uses the interaction-force information, and this
is true of every strategy, not a caveat limited to some of them.

Geometric tracking, INDI and feedback linearisation are different laws and are tuned separately.
After that tuning, every strategy's gains are frozen for the comparison flights so a later gap is
not mid-campaign retuning. The cost of per-strategy tuning is in scope for RQ4. The discussion
chapter reports where a frozen gain appears to have disadvantaged a strategy.

% -------------------------------------------------------------------------------------
\section{Summary}
\label{sec:ctrl-summary}

Three strategies (1--3: pure INDI, geometric+NN and FBL+NN), plus the uncompensated baseline
(Strategy~0) flown as the reference condition under all of them, each its own controller,
distinguished by what they know about the interaction force, when they know it, and where that
knowledge enters. All three are compared under the same experimental conditions: identical
trajectories, formations, hardware and separations, with each strategy's gains frozen after its own
tuning, not identical control structure.

At the time of writing, two strategies fly (the baseline and pure INDI), one is implemented and
unflown (geometric+NN), and one is not implemented (FBL+NN). Chapter~\ref{ch:setup} describes the experimental
design; the two- and multi-robot results chapters report what was flown.
